\documentclass{article}
% --------
% PACKAGES
% --------
\usepackage{texmaths-symbols}
\usepackage{refcount}
\usepackage{lastpage}
\usepackage{fancyhdr}
\usepackage[a5paper]{geometry}
\usepackage{tabularray}
\usepackage{enumitem}
\usepackage[tikz]{mdframed}
\usepackage{texmaths-cycle4-BO_2020}
% --------
% GEOMETRY
% --------
\geometry{
    hmargin=18pt,
    top=50pt,
    headheight=20pt,
    headsep=25pt,
    textheight=500pt,
    footskip=20pt
}
% -------------------
% HEADERS AND FOOTERS
% -------------------
\makeatletter
\def\level#1{\def\@level@name{#1}%
    \setcounter{totalseances}0
    \setcounter{sequence}0}
\fancypagestyle{sequencage}
    {
        \fancyhf{}
        \fancyhead[L]{\small Séquençage Cycle 4 \@level@name}
        \fancyhead[R]{\small D'après le BOEN n$^\circ$ 31 du 30 juillet 2020}
        \fancyfoot[C]{\thepage/\pageref{LastPage}}
    }

% --------
% COUNTERS
% --------
\newcounter{totalseances}
\newcounter{seancenumber}
\newcounter{sequence}
% ------
% MACROS
% ------
\def\storecounter#1#2{%
    \addtocounter{#1}{-1}\refstepcounter{#1}\label{#2}}
\def\seancenumberseq#1{%
    \getrefnumber{ref:seancenumber\@level@name#1}}
\long\def\seance#1{%
    \par\noindent
    \stepcounter{seancenumber}\stepcounter{totalseances}%
    {\bfseries Séance \theseancenumber : }#1
    }
\long\def\dontshowseance#1{\stepcounter{seancenumber}\stepcounter{totalseances}}
\def\sequence#1{%
    \begingroup
        \let\'\@firstofone
        \let\`\@firstofone
        \let\ \@underscore
        \xdef\sequence@key{#1}
    \endgroup
    \stepcounter{sequence}\vskip5pt
    \begin{mdframed}[roundcorner=5pt]
        \begin{tblr}{colspec={X[1,l]X[2,c]X[1,r]}}%{colspec={X[1,l]X[2,l]X[1,r]}}
            \bfseries Séquence \arabic{sequence}&
            \footnotesize#1&
            (\seancenumberseq\thesequence\ séances)
        \end{tblr}
    \end{mdframed}\vskip8pt
}
\def\endsequence{%
    \storecounter{seancenumber}{ref:seancenumber\@level@name\thesequence}%
    \setcounter{seancenumber}0%
}
\setlist{label=$-$}

\def\cnsandcmp#1#2{%
    {\itshape Connaissances}\vskip8pt
    \@for\ii:=#1\do{%
        \par\noindent$-$\hskip2ex\usecns[\ii]%
        \immediate\write\seqcanal{CNS\ii:\@level@name:\sequence@key:}%
    }\vskip5pt\noindent\nopagebreak
    {\itshape Compétences}\vskip8pt
    \@for\jj:=#2\do{%
    \par\noindent$-$\hskip2ex\usecmp[\jj]%
    % \immediate\write\seqcanal{CMP\jj:\@level@name:\sequence@key}%
    }\par\kern20pt\noindent
}

\def\showtotalseance{\vskip0pt plus 1fill\noindent
{\large Nombre total de séances en \@level@name: {\bfseries\thetotalseances}}
\vskip0pt plus 3fill}

\def\@underscore{\string_}

\def\sequence@key{undefined}

\makeatother
\def\hideseance#1{\stepcounter{seancenumber}\stepcounter{totalseances}}
% \let\seance\hideseance
\begin{document}

\newwrite\seqcanal
\immediate\openout\seqcanal=\jobname-structure.txt%

Les connaissances et compétences liées au thème E (algorithmique et
programmation) sont à travailler tout au long du cycle 4, en salle informatique
ou en exercice en classe, en lien avec d'autres notions.

\pagestyle{sequencage}

\level{5eme}

\sequence{Expressions numériques}
\cnsandcmp{A1h1,A1h2}{A1f1,A1f6,A1j1}
\seance{Vocabulaire}
\seance{Règles de priorités opératoires}
\seance{Traduction}
\seance{Version}
\seance{Résolution de problèmes}
\endsequence
\par\vskip10pt\noindent{\bfseries Automatismes : tracer figures usuelles à main levée}

\sequence{Constructibilité d'un triangle}
\cnsandcmp{D2b7,D2b4}{D2a*}
\seance{Sens géométrique de l'inégalité}
\seance{Figures à main levée qu'on ne peut pas tracer}
\seance{Application de l'inégalité}
\seance{Une seule inégalité suffit}
\seance{Résolution de problèmes}
\endsequence

\sequence{Découverte des nombres relatifs}
\cnsandcmp{A1a1,A1a2,A1a3,A1g1,A1h3}{A1a1,A1a4,A1c1,A1d1,A1j2}
\seance{Introduction aux nombres négatifs}
\seance{Comparer des nombres négatifs}
\seance{Notion d'opposé}
\seance{Repérage sur un plan}
\seance{Résolution de problèmes}
\seance{Additions}
\endsequence

\sequence{Opérations sur les fractions}
\cnsandcmp{A1b1,A2e*,A1h6,A1g2}{A1a2,A1b*,A1f4,A1j1,A1c2,A1d2}
\seance{Rappels : définition d'une fraction}
\seance{Activité : plusieurs fractions peuvent représenter un même nombre}
\seance{Modification d'une fraction / simplification}
\seance{Additions et soustractions de fractions de même dénominateur}
\seance{Additions de fractions de dénominateurs multiples}
\seance{Additions de fractions de dénominateurs quelconque}
\seance{Résolution de problèmes}
\endsequence

\sequence{Cas d'égalité des triangles}
\cnsandcmp{D2b5}{D2b*,D2d*}
\seance{Initiation à la démonstration}
\seance{CCC}
\seance{Exercices de démonstrations utilisant CCC}
\seance{CAC}
\seance{Exercices de démonstrations utilisant CAC}
\seance{ACA}
\seance{Exercices de démonstrations utilisant ACA}
\seance{Cas particuliers}
\seance{Résolution de problèmes}
\endsequence

\sequence{Calcul littéral}
\cnsandcmp{A3a1,A3a2,A3a3}{A3a3,A3b*}
\seance{Expression littérale}
\seance{Conventions d'écriture}
\seance{Evaluer une expression littérale}
\seance{Tester une égalité}
\seance{Réduction (uniquement des additions/soustrations)}
\seance{Réduction (uniquement des multiplications/divisions)}
\seance{Réduction (toutes opérations)}
\seance{Modélisation}
\endsequence

\sequence{Angles alternes-internes}
\cnsandcmp{D2a*}{D2d*,D2e*}
\seance{Égalité des angles opposés par le sommet}
\seance{Vocabulaire}
\seance{Sens direct}
\seance{Démonstrations (exercices)}
\seance{Réciproque}
\seance{Démonstrations (exercices)}
\endsequence

\sequence{Proportionnalité}
\cnsandcmp{B3a*,A1b2}{B3a*,B3e*}
\seance{Multiplications à trou}
\seance{Situation de proportionnalité / définition}
\seance{Méthodes additives / multiplicatives}
\seance{Coefficient de proportionnalité}
\seance{Représentation de la proportionnalité : tableau}
\seance{Résolution de problèmes}
\endsequence

\sequence{Somme des angles d'un triangle}
\cnsandcmp{D2b1}{C1a*,C1b*,C1c*,D2d*,D2e*}
\seance{Sommes des angles}
\seance{Démonstration de la propriété}
\seance{Retour sur ACA}
\endsequence

\sequence{Arithmétique}
\cnsandcmp{A2a*,A2b*,A2c*}{A2a*,A2c*,A2g*}
\seance{Division euclidienne}
\seance{Multiples et diviseurs}
\seance{Critères de divisibilité}
\seance{Lien avec les fractions}
\seance{Résolution de problème}
\endsequence

\sequence{Hauteurs et médiatrices d'un triangle}
\cnsandcmp{D2b2,D2b3}{D2a*}
\seance{Hauteurs}
\seance{Les hauteurs d'un triangle sont concourantes (activité)}
\seance{Médiatrices}
\seance{Les médiatrices d'un triangles sont concourantes (activité)}
\endsequence

\sequence{Additions et soustractions de nombres relatifs}
\cnsandcmp{A1h3,A1h4}{A1a1,A1f2,A1j1}
\seance{Additions de \guillemets{petits} nombre entiers : lien avec les dettes, température,~...}
\seance{Additions de nombres relatifs de même signe}
\seance{Additions de nombres relatifs signes}
\seance{Additions de plusieurs nombres relatifs}
\seance{Soustractions de nombres relatifs}
\seance{Enlever les parenthèses}
\seance{Enlever les parenthèses}
\endsequence

\sequence{Symétrie et centrale}
\cnsandcmp{C2a1}{C2c*,D2c1,D2d*,D2e*}
\seance{Rappels sur la symétrie axiale : construction}
\seance{Axes de symétrie des figures usuelles}
\seance{Symétrie centrale : centre de symétrie}
\seance{Définition de la symétrie centrale comme demi-tour}
\seance{Définition rigoureuse et construction}
\seance{Construction de figures complexes}
\endsequence

\sequence{Parallélogrammes}
\cnsandcmp{D2c*}{D2d*,D2e*}
\seance{Définition quadrilatère possédant un axe de symétrie}
\seance{Propriétés : angles opposés égaux et côtés opposés égaux}
\seance{Lien avec les rectangles / carrés}
\seance{Lien entre les diagonales et la nature du quadrilatère}
\endsequence

\showtotalseance\newpage

\level{4eme}
\sequence{Sommes\ et\ diff\'erences\ de\ nombres\ relatifs}
\cnsandcmp{A1h4}{A1f3,A1j2}
\seance{}
\seance{}
\seance{}
\seance{}
\endsequence
\sequence{Proportionnalit\'e}
\cnsandcmp{B3b*,B3c*,C1a*,B3c*}{A1a1,B3a*,B3b*,B3d*,B3e*}
\seance{}
\seance{}
\seance{}
\seance{}
\endsequence
\sequence{Produits\ et\ quotients\ de\ nombres\ relatifs}
\cnsandcmp{A1h5}{A1a1,A1b*,A1h2,A1j1}
\seance{}
\seance{}
\seance{}
\seance{}
\seance{}
\seance{}
\seance{}
\endsequence
\sequence{Translation}
\cnsandcmp{C2a2}{D2a*,D2c2,D2d*,D2e*}
\seance{}
\seance{}
\seance{}
\seance{}
\endsequence
\sequence{Sommes\ et\ diff\'erences\ de\ nombres\ rationnels}
\cnsandcmp{A1f*,A1b3,A2e*,A1h6}{A1a1,A1a2,A1f2,A2f*,A1d*,A1f5,A1j1}
\seance{}
\seance{}
\seance{}
\seance{}
\seance{}
\seance{}
\seance{}
\endsequence
\sequence{Rotation}
\cnsandcmp{C2a3}{D2a*,C2c*,D2c3,D2d*,D2e*}
\seance{}
\seance{}
\seance{}
\seance{}
\endsequence
\sequence{Produits\ et\ quotients\ de\ nombres\ rationnels}
\cnsandcmp{A1b4,A1h8,A1h9}{A1a2,A1b*,A1f5,A1j1}
\seance{}
\seance{}
\seance{}
\seance{}
\seance{}
\seance{}
\seance{}
\endsequence
\sequence{Calcul\ litt\'eral}
\cnsandcmp{A3a1,A3a2,A3a3}{A1a2,A3a3,A3b*}
\seance{}
\seance{}
\seance{}
\seance{}
\endsequence
\sequence{\'Equation}
\cnsandcmp{A3a4,A3a5}{A3c*,A3d1}
\seance{}
\seance{}
\seance{}
\seance{}
\endsequence
\sequence{Th\'eor\`eme\ de\ Pythagore}
\cnsandcmp{A1c*,A1d*,D2e1,C1b*}{A1b*,A1i*,D2d*,D2e*}
\seance{}
\seance{}
\seance{}
\seance{}
\seance{}
\seance{}
\seance{}
\endsequence
\sequence{Puissance\ d'un\ nombre}
\cnsandcmp{A1e*,A1i*,A1j*}{A1a3,A1b*,A1c3,A1h1,A1h2,A1j1,A1g*}
\seance{}
\seance{}
\seance{}
\seance{}
\seance{}
\seance{}
\seance{}
\endsequence
\sequence{R\'eciproque\ du\ th\'eor\`eme\ de\ Pythagore}
\cnsandcmp{D2e2}{D2d*,D2e*}
\seance{}
\seance{}
\seance{}
\seance{}
\endsequence
\sequence{Distributivit\'e}
\cnsandcmp{A3b1,A3b2}{A3a1,A3a2,A3b*}
\seance{}
\seance{}
\seance{}
\seance{}
\endsequence
\sequence{Statistiques}
\cnsandcmp{B1a*,B1b*}{B1a*,B1b*,B1d*,B1e*}
\seance{}
\seance{}
\seance{}
\seance{}
\endsequence
\sequence{Probabilit\'es}
\cnsandcmp{B2a*,B2b*,B2c*}{B2a*,B2b*}
\seance{}
\seance{}
\seance{}
\seance{}
\endsequence


\showtotalseance\newpage

\level{3eme}

\sequence{Arithm\'etique}
\cnsandcmp{A2d*}{A2b*,A2d*,A2e1,A2e2,A2g*}
\seance{Vocabulaire : diviseurs et multiples}
\seance{Définition d'un nombre premier et Crible d'Ératostène}
\seance{Décomposition d'un nombre en facteurs premiers : méthode 1}
\seance{Décomposition d'un nombre en facteurs premiers : méthode 2}
\seance{Exercice(s) type brevet}
\seance{Recherche du PGCD}
\seance{Exercice(s) type brevet}
\seance{Simplification de fractions}
\seance{Exercice(s) type brevet}
\endsequence

\sequence{Thal\`es\ (Partie\ 1)}
\cnsandcmp{D2b6,D2d1}{D2d*,D2e*}
\seance{Définition de deux triangle semblables}
\seance{Propriété des triangles semblables}
\seance{Théorème de Thalès : sens direct}
\seance{Exercices type brevet}
\seance{Théorème de Thalès : sens direct (configuration papillon)}
\seance{Exercices type brevet}
\endsequence

\sequence{Calcul\ litt\'eral}
\cnsandcmp{A3b1,A3b2,A3d*}{A3a1,A3a2,A3b*}
\seance{Réduction d'une expression littérale}
\seance{Découverte de la propriété de la distributivité simple}
\seance{Distributivité simple : développer dans des cas simples}
\seance{Distributivité simple : développer puis réduire}
\seance{Scratch + exercices type brevet}
\seance{Distributivité simple : factoriser}
\seance{Distributivité simple : factoriser (lien avec le regroupement de termes)}
\seance{Distributivité double : découverte de la propriété + démonstration}
\seance{Distributivité double : développer}
\seance{exercices types brevet}
\seance{Identités remarquables}
\seance{Identité remarquable : factoriser $a^2-b^2$}
\endsequence

\sequence{Statistiques}
\cnsandcmp{B1a*,B1b*,B1c*}{B1a*,B1b*,B1c*,B1d*,B1e*}
\seance{Définition d'une série statistique et calcul de moyenne (cas simple)}
\seance{Calcul de médiane (cas simple) et étendue}
\seance{Utilité de la notion d'effectifs : tableau d'effectifs}
\seance{Fréquence + exercice type brevet}
\seance{Calcul de la moyenne pondérée}
\seance{Calcul de la médiane avec les effectifs}
\seance{Exercices types brevet}
\seance{Lecture de diagramme (tableau, bâtons)}
\seance{Lecture de diagramme (histogramme)}
\seance{Lecture de diagramme (circulaire)}
\seance{Exercices type brevet}
\endsequence

\sequence{\'Equations}
\cnsandcmp{A3c*}{A3c*,A3d1,A3d2}
\seance{Rappels sur les équations : vocabulaire}
\seance{Résolutions d'équations par recherche des solutions}
\seance{Équations du type $ax+b=c$ : technique}
\seance{Équations du type $ax+b=c$ : mise en équation}
\seance{Équations du type $ax+b=cx+d$ : technique}
\seance{Équations du type $ax+b=cx+d$ : mise en équation}
\seance{Équation produit nul : technique}
\seance{Mise en équations}
\seance{Equation carré constant}
\seance{Mise en équations}
\endsequence

\sequence{Notion\ de\ fonction}
\cnsandcmp{B4a*,B4b*,B4c*}{B4a*,B4b*,B4c*,B4d*,B4f*}
\seance{Découverte de la notion de fonction, vocabulaire}
\seance{Exprimer une grandeur en fonction de $x$}
\seance{Utiliser le vocabulaire (antécédent, image)}
\seance{Tableau de valeur}
\seance{Exercices type brevet}
\seance{Représentation graphique}
\seance{Exercices type brevet}
\endsequence

\sequence{Trigonom\'etrie}
\cnsandcmp{D2f*}{D2d*,D2e*}
\seance{Découverte des fonction cos et sin: côtés adj/opp dans un triangle d'hypoténuse 1}
\seance{Découverte des trois formules trigonométriques: deux manières de les calculer}
\seance{Calculer un angle en connaissant les lignes trigonométriques}
\seance{Calculer une longueur en connaissant les angles}
\endsequence

\sequence{Thal\`es\ (Partie\ 2)}
\cnsandcmp{D2d2}{D2e*}
\endsequence

\sequence{Homoth\'etie}
\cnsandcmp{C2a4}{C2a*,C2b*,C2c*,C2d*,D2c4}
\endsequence

\sequence{Fonctions\ lin\'eaires\ et\ affine}
\cnsandcmp{B4d*}{B3a*,B3d*,B4e*}
\endsequence

\sequence{Probabilit\'es}
\cnsandcmp{B2a*,B2b*,B2c*}{B2a*,B2b*,B2c*,B2d*}
\endsequence

\sequence{Repr\'esenter\ l'espace}
\cnsandcmp{D1a*,D1b*,C1c1,C1c2,C1c3,C1c4,C1c5,C1d*}{D1a*,D1b*,D1c*,D1d*,C1a*}
\endsequence

\showtotalseance



\immediate\closeout\seqcanal

\newpage
\makeatletter
\foreach\themefor in{1,...,\fourthcycle@total[themes]}%
{%
    \newpage
    \def\themecode{\nbtoAlph\themefor}
    {\par\centering\LARGE\bfseries Thème \themecode\ : \fourthcycle@theme[\themecode]\centering
        \par}\vskip8pt
    \foreach\attendufor in{1,...,\fourthcycle@total[\nbtoAlph\themefor]}
    {%
        \def\attenducode{\nbtoAlph\themefor\attendufor}
        \begin{mdframed}\itshape\fourthcycle@attendu[\attenducode]\hbadness=500\end{mdframed}
        {\bfseries Connaissances}\vskip8pt
        \foreach\cnsfor in{1,...,\fourthcycle@total[\nbtoAlph\themefor\attendufor CNS]}
            {\def\cnscode{\nbtoAlph\themefor\attendufor\nbtoalph\cnsfor}%
            {\tt\cnscode\hskip10pt}
            \thetoks{CNS\cnscode*occ}
            \vskip8pt
            \begingroup\setlength{\leftskip}{10pt}
            \ifnum\subcnstotal>0
                \foreach\subcnsfor in{1,...,\subcnstotal}{
                    {\itshape\subcnsfor)}\hskip5pt
                    \thetoks{CNS\nbtoAlph\themefor\attendufor
                    \nbtoalph\cnsfor\subcnsfor occ}\vskip5pt
                }%
            \fi\par\endgroup
        }%
        {\bfseries Compétences}\vskip8pt
        \foreach\cmpfor in{1,...,\fourthcycle@total[\nbtoAlph\themefor\attendufor CMP]}
        {%
            \def\cmpcode{\nbtoAlph\themefor\attendufor\nbtoalph\cmpfor}%
            {\tt\cmpcode\hskip10pt}
            \thetoks{CMP\cmpcode*occ}
            \vskip8pt
            \begingroup\setlength{\leftskip}{10pt}
            \ifnum\subcmptotal>0
                \foreach\subcmpfor in{1,...,\subcmptotal}
                {%
                {\itshape\subcmpfor)}\hskip5pt
                \thetoks{CMP\nbtoAlph\themefor\attendufor
                \nbtoalph\cmpfor\subcmpfor occ}\vskip5pt
                }%
            \fi\par\endgroup
        }%
    }%
}%

\end{document}